\section{Factorization}
\label{section-factoring}

\noindent
Here are some notions and relations between them that are typically taught
in a first year course on algebra at the undergraduate level.

\begin{definition}
\label{definition-irreducible-prime-element}
Let $R$ be a domain.
\begin{enumerate}
\item Elements $x, y \in R$ are called {\it associates} if
there exists a unit $u \in R^*$ such that $x = uy$.
\item An element $x \in R$ is called {\it irreducible}
if it is nonzero, not a unit and whenever $x = yz$, $y, z \in R$,
then $y$ is either a unit or an associate of $x$.
\item An element $x \in R$ is called {\it prime} if the ideal
generated by $x$ is a prime ideal.
\end{enumerate}
\end{definition}

\begin{lemma}
\label{lemma-easy-divisibility}
Let $R$ be a domain. Let $x, y \in R$.
Then $x$, $y$ are associates if and only if $(x) = (y)$.
\end{lemma}

\begin{proof}
If $x = uy$ for some unit $u \in R$, then $(x) \subset (y)$ and
$y = u^{-1}x$ so also $(y) \subset (x)$. Conversely, suppose that
$(x) = (y)$. Then $x = fy$ and $y = gx$ for some $f, g \in A$.
Then $x = fg x$ and since $R$ is a domain $fg = 1$. Thus
$x$ and $y$ are associates.
\end{proof}

\begin{lemma}
\label{lemma-factorization-exists}
Let $R$ be a domain. Consider the following conditions:
\begin{enumerate}
\item The ring $R$ satisfies the ascending chain condition for
principal ideals.
\item Every nonzero, nonunit element $a \in R$
has a factorization $a = b_1 \ldots b_k$
with each $b_i$ an irreducible element of $R$.
\end{enumerate}
Then (1) implies (2).
\end{lemma}

\begin{proof}
Let $x$ be a nonzero element, not a unit, which does not have a
factorization into
irreducibles. Set $x_1 = x$. We can write $x = yz$ where neither
$y$ nor $z$ is irreducible or a unit.
Then either $y$ does not have a factorization
into irreducibles, in which case we set $x_2 = y$, or $z$ does not
have a factorization into irreducibles, in which case we set $x_2 = z$.
Continuing in this fashion we find a sequence
$$
\ldots | x_3 | x_2 | x_1
$$
of elements of $R$ with $x_n/x_{n + 1}$ not a unit.
This gives a strictly increasing sequence of principal ideals
$(x_1) \subset (x_2) \subset (x_3) \subset \ldots$ thereby finishing the proof.
\end{proof}

\begin{definition}
\label{definition-UFD}
A {\it unique factorization domain}, abbreviated {\it UFD},
is a domain $R$ such that
if $x \in R$ is a nonzero, nonunit, then $x$ has a factorization
into irreducibles, and if
$$
x = a_1 \ldots a_m = b_1 \ldots b_n
$$
are factorizations into irreducibles then $n = m$ and
there exists a permutation $\sigma : \{1, \ldots, n\} \to \{1, \ldots, n\}$
such that $a_i$ and $b_{\sigma(i)}$ are associates.
\end{definition}

\begin{lemma}
\label{lemma-characterize-UFD}
Let $R$ be a domain. Assume every nonzero, nonunit factors into
irreducibles. Then $R$ is a UFD if and only if every irreducible
element is prime.
\end{lemma}

\begin{proof}
Assume $R$ is a UFD and let $x \in R$ be an irreducible element.
Say $ab \in (x)$, i.e., $ab = cx$. Choose factorizations
$a = a_1 \ldots a_n$, $b = b_1 \ldots b_m$, and $c = c_1 \ldots c_r$.
By uniqueness of the factorization
$$
a_1 \ldots a_n b_1 \ldots b_m = c_1 \ldots c_r x
$$
we find that $x$ is an associate of one of the elements
$a_1, \ldots, b_m$. In other words, either $a \in (x)$ or $b \in (x)$
and we conclude that $x$ is prime.

\medskip\noindent
Assume every irreducible element is prime. We have to prove that
factorization into irreducibles is unique up to permutation and
taking associates. Say $a_1 \ldots a_m = b_1 \ldots b_n$ with $a_i$
and $b_j$ irreducible. Since $a_1$ is prime, we see that
$b_j \in (a_1)$ for some $j$. After renumbering we may assume
$b_1 \in (a_1)$. Then $b_1 = a_1 u$ and since $b_1$ is irreducible
we see that $u$ is a unit. Hence $a_1$ and $b_1$ are associates
and $a_2 \ldots a_n = ub_2\ldots b_m$. By induction on $n + m$
we see that $n = m$ and $a_i$ associate to $b_{\sigma(i)}$ for
$i = 2, \ldots, n$ as desired.
\end{proof}

\begin{lemma}
\label{lemma-characterize-UFD-height-1}
Let $R$ be a Noetherian domain. Then $R$ is a UFD if and only if every
height $1$ prime ideal is principal.
\end{lemma}

\begin{proof}
Assume $R$ is a UFD and let $\mathfrak p$ be a height 1 prime ideal.
Take $x \in \mathfrak p$ nonzero and let $x = a_1 \ldots a_n$
be a factorization into irreducibles.
Since $\mathfrak p$ is prime we see that $a_i \in \mathfrak p$
for some $i$. By Lemma \ref{lemma-characterize-UFD} the ideal
$(a_i)$ is prime. Since $\mathfrak p$ has height $1$ we conclude
that $(a_i) = \mathfrak p$.

\medskip\noindent
Assume every height $1$ prime is principal. Since $R$ is Noetherian
every nonzero nonunit element $x$ has a factorization into irreducibles,
see Lemma \ref{lemma-factorization-exists}. It suffices to prove that
an irreducible element $x$ is prime, see Lemma \ref{lemma-characterize-UFD}.
Let $(x) \subset \mathfrak p$ be a prime minimal over $(x)$. Then
$\mathfrak p$ has height $1$ by Lemma \ref{lemma-minimal-over-1}.
By assumption $\mathfrak p = (y)$. Hence $x = yz$ and $z$ is a unit
as $x$ is irreducible. Thus $(x) = (y)$ and we see that $x$ is prime.
\end{proof}

\begin{lemma}[Nagata's criterion for factoriality]
\label{lemma-invert-prime-elements}
\begin{reference}
\cite[Lemma 2]{Nagata-UFD}
\end{reference}
Let $A$ be a domain. Let $S \subset A$ be a multiplicative subset
generated by prime elements. Let $x \in A$ be irreducible. Then
\begin{enumerate}
\item the image of $x$ in $S^{-1}A$ is irreducible or a unit, and
\item $x$ is prime if and only if the image of $x$ in $S^{-1}A$ is
a prime element or a unit in $S^{-1}A$.
\end{enumerate}
Moreover, then $A$ is a UFD if and only if every nonzero nonunit
element of $A$ has a factorization into irreducibles and $S^{-1}A$ is a UFD.
\end{lemma}

\begin{proof}
Say $x = \alpha \beta$ for $\alpha, \beta \in S^{-1}A$. Then
$\alpha = a/s$ and $\beta = b/s'$ for $a, b \in A$, $s, s' \in S$.
Thus we get $ss'x = ab$. By assumption we can write
$ss' = p_1 \ldots p_r$ for some prime elements $p_i$.
For each $i$ the element $p_i$ divides either $a$ or $b$.
Dividing we find a factorization $x = a' b'$ and
$a = s'' a'$, $b = s''' b'$ for some $s'', s''' \in S$.
As $x$ is irreducible, either $a'$ or $b'$ is a unit.
Tracing back we find that either $\alpha$ or $\beta$ is a unit.
This proves (1).

\medskip\noindent
Suppose $x$ is prime. Then $A/(x)$ is a domain. Hence
$S^{-1}A/xS^{-1}A = S^{-1}(A/(x))$ is a domain or zero.
Thus $x$ maps to a prime element or a unit.

\medskip\noindent
Suppose that the image of $x$ in $S^{-1}A$ is a unit.
Then $y x = s$ for some $s \in S$ and $y \in A$. By assumption
$s = p_1 \ldots p_r$ with $p_i$ a prime element. For each $i$ either
$p_i$ divides $y$ or $p_i$ divides $x$. In the second case
$p_i$ and $x$ are associates (as $x$ is irreducible) and we are done.
But if the first case happens for all $i = 1, \ldots, r$, then
$x$ is a unit which is a contradiction.

\medskip\noindent
Suppose that the image of $x$ in $S^{-1}A$ is a prime element.
Assume $a, b \in A$ and $ab \in (x)$. Then $sa = xy$ or
$sb = xy$ for some $s \in S$ and $y \in A$. Say the first
case happens. By assumption
$s = p_1 \ldots p_r$ with $p_i$ a prime element. For each $i$ either
$p_i$ divides $y$ or $p_i$ divides $x$. In the second case
$p_i$ and $x$ are associates (as $x$ is irreducible) and we are done.
If the first case happens for all $i = 1, \ldots, r$, then
$a \in (x)$ as desired. This completes the proof of (2).

\medskip\noindent
The final statement of the lemma follows from (1) and (2)
and Lemma \ref{lemma-characterize-UFD}.
\end{proof}

\begin{lemma}
\label{lemma-UFD-ascending-chain-condition-principal-ideals}
A UFD satisfies the ascending chain condition for principal
ideals.
\end{lemma}

\begin{proof}
Consider an ascending chain $(a_1) \subset (a_2) \subset (a_3) \subset \ldots$
of principal ideals in $R$. Write $a_1 = p_1^{e_1} \ldots p_r^{e_r}$
with $p_i$ prime. Then we see that $a_n$ is an associate of
$p_1^{c_1} \ldots p_r^{c_r}$ for some $0 \leq c_i \leq e_i$.
Since there are only finitely many possibilities we conclude.
\end{proof}

\begin{lemma}
\label{lemma-factoring-in-polynomial}
Let $R$ be a domain. Assume $R$ has the ascending chain condition
for principal ideals. Then the same property holds for a polynomial
ring over $R$.
\end{lemma}

\begin{proof}
Consider an ascending chain $(f_1) \subset (f_2) \subset (f_3) \subset \ldots$
of principal ideals in $R[x]$. Since $f_{n + 1}$ divides $f_n$ we see
that the degrees decrease in the sequence. Thus $f_n$
has fixed degree $d \geq 0$ for all $n \gg 0$. Let $a_n$ be the
leading coefficient of $f_n$. The condition $f_n \in (f_{n + 1})$
implies that $a_{n + 1}$ divides $a_n$ for all $n$.
By our assumption on $R$ we see that $a_{n + 1}$ and $a_n$
are associates for all $n$ large enough (Lemma \ref{lemma-easy-divisibility}).
Thus for large $n$ we see that $f_n = u f_{n + 1}$ where
$u \in R$ (for reasons of degree) is a unit (as $a_n$ and $a_{n + 1}$
are associates).
\end{proof}

\begin{lemma}
\label{lemma-polynomial-ring-UFD}
A polynomial ring over a UFD is a UFD. In particular, if $k$ is a field,
then $k[x_1, \ldots, x_n]$ is a UFD.
\end{lemma}

\begin{proof}
Let $R$ be a UFD. Then $R$ satisfies the ascending chain condition for
principal ideals
(Lemma \ref{lemma-UFD-ascending-chain-condition-principal-ideals}),
hence $R[x]$ satisfies the ascending chain condition for principal
ideals (Lemma \ref{lemma-factoring-in-polynomial}), and hence every
element of $R[x]$ has a factorization into irreducibles
(Lemma \ref{lemma-factorization-exists}).
Let $S \subset R$ be the multiplicative subset
generated by prime elements. Since every nonunit of $R$ is a product
of prime elements we see that $K = S^{-1}R$ is the fraction field
of $R$. Observe that every prime element of $R$ maps to a prime element
of $R[x]$ and that $S^{-1}(R[x]) = S^{-1}R[x] = K[x]$ is a
UFD (and even a PID). Thus we may apply
Lemma \ref{lemma-invert-prime-elements} to conclude.
\end{proof}

\begin{lemma}
\label{lemma-UFD-normal}
A unique factorization domain is normal.
\end{lemma}

\begin{proof}
Let $R$ be a UFD. Let $x$ be an element of the fraction field of
$R$ which is integral over $R$. Say $x^d - a_1 x^{d - 1} - \ldots - a_d = 0$
with $a_i \in R$. We can write
$x = u p_1^{e_1} \ldots p_r^{e_r}$ with $u$ a unit, $e_i \in \mathbf{Z}$,
and $p_1, \ldots, p_r$ irreducible elements which are not associates.
To prove the lemma we have to show $e_i \geq 0$. If not, say $e_1 < 0$,
then for $N \gg 0$ we get
$$
u^d p_2^{de_2 + N} \ldots p_r^{de_r + N} =
p_1^{-de_1}p_2^N \ldots p_r^N(
\sum\nolimits_{i = 1, \ldots, d} a_i x^{d - i} ) \in (p_1)
$$
which contradicts uniqueness of factorization in $R$.
\end{proof}

\begin{definition}
\label{definition-PID}
A {\it principal ideal domain}, abbreviated {\it PID},
is a domain $R$ such that every ideal is a principal ideal.
\end{definition}

\begin{lemma}
\label{lemma-PID-UFD}
A principal ideal domain is a unique factorization domain.
\end{lemma}

\begin{proof}
As a PID is Noetherian this follows from
Lemma \ref{lemma-characterize-UFD-height-1}.
\end{proof}

\begin{definition}
\label{definition-dedekind-domain}
A {\it Dedekind domain} is a domain $R$ such that every
nonzero ideal $I \subset R$ can be written as a product
$$
I = \mathfrak p_1 \ldots \mathfrak p_r
$$
of nonzero prime ideals uniquely up to permutation of the $\mathfrak p_i$.
\end{definition}

\begin{lemma}
\label{lemma-PID-dedekind}
A PID is a Dedekind domain.
\end{lemma}

\begin{proof}
Let $R$ be a PID. Since every nonzero ideal of $R$ is principal,
and $R$ is a UFD (Lemma \ref{lemma-PID-UFD}), this follows from
the fact that every irreducible element in $R$ is prime
(Lemma \ref{lemma-characterize-UFD})
so that factorizations of elements turn into factorizations into primes.
\end{proof}

\begin{lemma}
\label{lemma-product-ideals-principal}
\begin{slogan}
A product of ideals is an invertible module iff both factors are.
\end{slogan}
Let $A$ be a ring. Let $I$ and $J$ be nonzero ideals of $A$
such that $IJ = (f)$ for some nonzerodivisor $f \in A$. Then $I$ and $J$ are
finitely generated ideals and finitely locally free of rank $1$ as $A$-modules.
\end{lemma}

\begin{proof}
It suffices to show that $I$ and $J$ are finite locally free $A$-modules
of rank $1$, see Lemma \ref{lemma-finite-projective}.
To do this, write $f = \sum_{i = 1, \ldots, n} x_i y_i$ with $x_i \in I$
and $y_i \in J$. We can
also write $x_i y_i = a_i f$ for some $a_i \in A$.
Since $f$ is a nonzerodivisor we see that $\sum a_i = 1$.
Thus it suffices to show that each $I_{a_i}$ and
$J_{a_i}$ is free of rank $1$ over $A_{a_i}$. After replacing $A$ by
$A_{a_i}$ we conclude that $f = xy$ for some $x \in I$ and $y \in J$.
Note that both $x$ and $y$ are nonzerodivisors. We claim that
$I = (x)$ and $J = (y)$ which finishes the proof. Namely, if $x' \in I$,
then $x'y = af = axy$ for some $a \in A$. Hence $x' = ax$ and we win.
\end{proof}

\begin{lemma}
\label{lemma-characterize-Dedekind}
Let $R$ be a ring. The following are equivalent
\begin{enumerate}
\item $R$ is a Dedekind domain,
\item $R$ is a Noetherian domain and for every nonzero maximal ideal
$\mathfrak m$ the local ring $R_{\mathfrak m}$ is a discrete valuation ring, and
\item $R$ is a Noetherian, normal domain, and $\dim(R) \leq 1$.
\end{enumerate}
\end{lemma}

\begin{proof}
Assume (1). The argument is nontrivial because we did not assume that $R$
was Noetherian in our definition of a Dedekind domain. Let
$\mathfrak p \subset R$ be a prime ideal. Observe that
$\mathfrak p \not = \mathfrak p^2$ by uniqueness
of the factorizations in the definition. Pick $x \in \mathfrak p$
with $x \not \in \mathfrak p^2$. Let $y \in \mathfrak p$ be a
second element (for example $y = 0$).
Write $(x, y) = \mathfrak p_1 \ldots \mathfrak p_r$.
Since $(x, y) \subset \mathfrak p$ at least one of the primes
$\mathfrak p_i$ is contained in $\mathfrak p$.
But as $x \not \in \mathfrak p^2$ there is at most one.
Thus exactly one of $\mathfrak p_1, \ldots, \mathfrak p_r$ is contained
in $\mathfrak p$, say $\mathfrak p_1 \subset \mathfrak p$.
We conclude that $(x, y)R_\mathfrak p = \mathfrak p_1R_\mathfrak p$
is prime for every choice of $y$. We claim that
$(x)R_\mathfrak p = \mathfrak pR_\mathfrak p$. Namely,
pick $y \in \mathfrak p$. By the above applied with $y^2$ we see
that $(x, y^2)R_\mathfrak p$ is prime.
Hence $y \in (x, y^2)R_\mathfrak p$, i.e., $y = ax + by^2$ in
$R_\mathfrak p$. Thus $(1 - by)y = ax \in (x)R_\mathfrak p$, i.e.,
$y \in (x)R_\mathfrak p$ as desired.

\medskip\noindent
Writing $(x) = \mathfrak p_1 \ldots \mathfrak p_r$ anew with
$\mathfrak p_1 \subset \mathfrak p$ we conclude that
$\mathfrak p_1 R_\mathfrak p = \mathfrak p R_\mathfrak p$, i.e.,
$\mathfrak p_1 = \mathfrak p$. Moreover, $\mathfrak p_1 = \mathfrak p$ is
a finitely generated ideal of $R$ by
Lemma \ref{lemma-product-ideals-principal}.
We conclude that $R$ is Noetherian by Lemma \ref{lemma-cohen}.
Moreover, it follows that $R_\mathfrak m$ is a discrete
valuation ring for every prime ideal $\mathfrak p$, see
Lemma \ref{lemma-characterize-dvr}.

\medskip\noindent
The equivalence of (2) and (3) follows from
Lemmas \ref{lemma-normality-is-local} and
\ref{lemma-characterize-dvr}. Assume (2) and (3) are satisfied.
Let $I \subset R$ be an ideal. We will construct a factorization
of $I$. If $I$ is prime, then there is nothing to prove.
If not, pick $I \subset \mathfrak p$ with $\mathfrak p \subset R$
maximal. Let $J = \{x \in R \mid x \mathfrak p \subset I\}$.
We claim $J \mathfrak p = I$. It suffices to check this after localization
at the maximal ideals $\mathfrak m$ of $R$ (the formation of $J$ commutes with
localization and we use Lemma \ref{lemma-characterize-zero-local}).
Then either $\mathfrak p R_\mathfrak m = R_\mathfrak m$ and the result
is clear, or $\mathfrak p R_\mathfrak m = \mathfrak m R_\mathfrak m$.
In the last case $\mathfrak p R_\mathfrak m = (\pi)$ and the case where
$\mathfrak p$ is principal is immediate. By Noetherian induction the
ideal $J$ has a factorization and we obtain the desired factorization
of $I$. We omit the proof of uniqueness of the factorization.
\end{proof}

\noindent
The following is a variant of the Krull-Akizuki lemma.

\begin{lemma}
\label{lemma-integral-closure-Dedekind}
Let $A$ be a Noetherian domain of dimension $1$ with fraction field $K$.
Let $L/K$ be a finite extension. Let $B$ be the
integral closure of $A$ in $L$. Then $B$ is a Dedekind domain and
$\Spec(B) \to \Spec(A)$ is surjective, has finite fibres, and
induces finite residue field extensions.
\end{lemma}

\begin{proof}
By Krull-Akizuki (Lemma \ref{lemma-krull-akizuki})
the ring $B$ is Noetherian. By Lemma \ref{lemma-integral-sub-dim-equal}
$\dim(B) = 1$. Thus $B$ is a Dedekind domain by
Lemma \ref{lemma-characterize-Dedekind}.
Surjectivity of the map on spectra follows from
Lemma \ref{lemma-integral-overring-surjective}.
The last two statements follow from
Lemma \ref{lemma-finite-extension-residue-fields-dimension-1}.
\end{proof}







